\section{Fault Analysis and Reproducibility}
The fault matrix is more informative than a single headline reduction. Under no fault, fixed richer, static escalation, and decision driven all resolve the two attack scenario families and benign alternatives according to the declared rules. Under the denial fault, the distinguishing \texttt{failure\_kind} and \texttt{service\_errors} entries cannot be released. The 40 abnormal cases remain unresolved for decision driven; the ten normal cases have no candidate and remain \texttt{none}. Under a two-unit snapshot cap, the initial signal consumes one unit and there is insufficient space for all required predicates. Again 40 abnormal cases are unresolved. Under source-spread coarsening, the unknown category prevents confirmation of the two attack scenarios while some benign alternatives still contradict without it, yielding 20 unresolved cases. These patterns are consequences of the fixed rules and policy, not estimates of how frequently faults occur in production.

\begin{table}[t]
\centering\caption{Three disclosure strategies under each fault; cells are TP/FP/U across 50 cases.}
\label{tab:matrix}\scriptsize
\begin{tabular}{lccc}\hline
Fault & Richer & Static & Decision \\\hline
None &20/0/0&20/0/0&20/0/0\\
Deny &0/0/40&0/0/40&0/0/40\\
Coarsen &0/0/20&0/0/20&0/0/20\\
Budget &0/0/40&0/0/40&0/0/40\\\hline
\end{tabular}
\end{table}

The same 50 scenario--seed pairs are replayed for every strategy and fault. This pairing permits a direct within-harness comparison of totals but does not justify a confidence interval over unrelated organizations. The synthetic generator creates 10-second windows with 30 requests/s for normal, 200 requests/s for authentication attack and provider outage, and 800 requests/s for traffic attack and flash crowd. Authentication failures are marked with 401 status in the attack and 503 status in the outage. Traffic attacks inject elevated 503 responses. The source spread threshold is ten distinct addresses. These settings deliberately create two tempting benign alternatives for a coarse alert. They are disclosed here so a reader can see that the evaluation is a protocol test shaped around the implemented rules.

The result files are \texttt{research/results/disclosure.csv} and \texttt{research/results/summary.json}. The executable harness is \texttt{research/evaluate\_disclosure.py}; the case manifest and digest are \texttt{research/cases.json} and \texttt{research/cases.sha256}. The evaluator refuses to run if the manifest digest differs or the output CSV already exists, protecting an existing measurement from accidental overwrite. Each run creates separate organization ledger paths by fault, seed, scenario, and strategy. The summary script aggregates the recorded CSV rather than regenerating results. Tests cover API contracts, ledger behavior, durable recovery, and evidence lifecycle; integration checks also exercised migrations, backup/restore, frontend build, and a three-organization local demo. The paper's numerical claims should be rechecked against those files whenever the implementation changes.

\subsection{Interpretation of cost and transport}
The 18.2\% policy-unit reduction and 13.2\% serialized-byte reduction in the no-fault slice are different quantities. An approved \texttt{source\_spread} costs two units but its JSON body need not be twice as large as another predicate. Conversely, a zero-unit coarse response still consumes bytes and can reveal that the owner was queried. The measured payload also counts each request, which explains why fewer releases do not imply an equal percentage byte reduction. The static escalation strategy has the same units and bytes as fixed richer but a different round count because its schedule sends requests sequentially. No bandwidth measurement over TLS or WAN was recorded. A deployment claim about transfer cost therefore needs a network trace, compression and batching policy, and realistic retransmission rates.

\subsection{Sensitivity to correlation assumptions}
The harness uses three nodes and a three-vote outcome. If quorum were reduced to two, one faulty or colluding pair could dominate; if increased above three, these cases could never resolve. If the temporal bucket were wider, unrelated events could be combined; if narrower, concurrent observations might split. The manuscript reports configured defaults rather than claiming they are optimized. A deployment should measure clock skew, base rates, organization overlap, and correlated detector errors before selecting quorum or a time window. Provenance labels are crucial because three feeds forwarding one sensor are not three independent observations. The current tests verify deduplication logic, while empirical independence of real organizations remains unmeasured.

\subsection{Accounting under the fault interventions}
The cost matrix in Table~\ref{tab:units} separates a stopping-rule benefit from a policy constraint. Without faults, decision driven spends fewer units than fixed richer because it does not request the third predicate for two benign alternatives. Under denial, it spends half as many units as fixed richer because it stops at the first denied decisive entry for each affected source. The saved units in this condition have a different meaning from no-fault savings: the unresolved count rises to 40, so the coordinator has paid less partly because it could not finish. The coarsening condition gives the same total policy units for the three richer strategies, even though decision driven sends fewer bytes than static escalation. The budget condition binds before the full catalogue can be released, so all three richer strategies spend 240 units and leave 40 abnormal cases unresolved. These interventions demonstrate why a paper should report utility and exposure together.

\begin{table}[t]
\centering\caption{Aggregate policy units per 50-case fault slice.}
\label{tab:units}\scriptsize
\begin{tabular}{lrrrr}\hline
Strategy & None & Deny & Coarsen & Budget \\\hline
Local only &0&0&0&0\\
Fixed minimal &120&120&120&120\\
Fixed richer &660&480&420&240\\
Static escalation &660&480&420&240\\
Decision driven &540&240&420&240\\\hline
\end{tabular}
\end{table}

The fixed-minimal line is constant because it sends only the initial signal, which is allowed in these synthetic policies. Its apparent robustness to faults is misleading: it also produces the same 20 false positives in every slice and never tests the distinguishing predicates. Local-only incurs no sharing units, but it makes the same 20 false-positive support decisions and provides no cross-organization corroboration. These baselines bound two extremes: no cross-boundary evidence and initial-signal-only sharing. They are useful controls, not realistic final incident adjudicators. An analyst could manually resolve their alerts with data that the harness does not model.

The denial policy is especially instructive. In authentication, the first cheap failure-rate category can be approved but the later failure-kind category is refused. In traffic, request rate is approved but service errors are refused. The coordinator learns that an anomaly exists and learns one additional category, yet it cannot determine support or contradiction for either family. A naive implementation might count ``no contradicting response'' as support. The explicit unresolved state prevents that inference. It also has a cost: the organization has already released some information without delivering a final classification. If denials are frequent, an owner might prefer policy-aware negotiation before any release, such as advertising which catalogue entries it permits. That would reduce wasted partial disclosure but reveal policy configuration itself. We did not implement or evaluate such negotiation.

Coarsening creates a different ambiguity. The owner sends an explicit \texttt{unknown} source-spread value, so the coordinator knows the field was handled, yet the value cannot satisfy wide or narrow. Attack cases need wide spread to support and remain unresolved. Some benign alternatives can still contradict through failure kind or service-error values, so the unresolved count is 20 rather than 40. A coarsened response is therefore more informative about process than a denial but can be equally nondecisive for a particular rule. The implemented cost assigns it zero units because no fine value is released; whether even the existence of the coarse response deserves a charge is a policy choice for a real deployment.

\subsection{Reproducing a result without overstating it}
The case manifest is checked by digest before the evaluator writes a CSV. The evaluator refuses an existing output path, requiring a deliberate new result directory or a conscious removal of a prior artifact. This prevents accidental silent replacement of numbers cited by the manuscripts. Each output row records its scenario, seed, strategy, fault, ground-truth attack flag, outcome, unit cost, refinement count, serialized bytes, rounds, and local elapsed time. A reader can independently group rows by strategy and fault and recompute every table in this section. The summary JSON is a convenience view, while the row-level CSV preserves diagnostics when a total changes.

Reproducibility is not the same as independence. The pseudorandom seeds fix event realizations within a generator. They do not sample five new scenario types or three unrelated organizations. The normal and abnormal rates and status-code mappings are explicit in the simulator. If a researcher changes the detector threshold, catalogue order, or source rule, the protocol totals may change even with the same seeds. The scientific object to reproduce is the complete combination of case manifest, code revision, environment, and row-level output. This is why the manuscript confines its quantitative claim to the tested implementation and does not claim a statistically significant privacy or detection improvement.

There is also a construct-validity issue in the term \emph{false positive}. In this experiment it means a benign-labelled synthetic scenario received a \texttt{supported} protocol outcome. That is not necessarily a false security alert in operational terms: a flash crowd may deserve attention, and a provider outage is itself an incident. The term is used to compare the declared attack-labelled task. A deployment should report separate labels for malicious activity, service health, and analyst actionability. A single binary attack label cannot capture all three. Likewise, \emph{true positive} here means the rule supports a scenario authored as an attack; it does not verify that the detector would identify an unseen attack in live traffic.

\subsection{Potential adaptations}
There are plausible next policies beyond the two fixed schedules and the greedy selector. An owner could provide a signed capability map describing entries it may release at each coarseness, allowing the coordinator to avoid doomed requests. A coordinator could use a prior over benign alternatives to select a more informative predicate first. Or peers could compute a joint threshold using cryptographic techniques, reducing explicit categorical release. Each option shifts a different trade-off: capability maps leak policy, learned priors may be brittle or biased, and cryptographic protocols add computational and key-management cost. The present method deliberately chooses a deterministic finite baseline before those complexities. It gives researchers a concrete reference point for comparing them rather than claiming to dominate all alternatives.

\section{Authorization as a Time-Varying Input}
Permission is evaluated at two times: when the snapshot is constructed for generation and when a report is retrieved. The first check prevents the model from seeing evidence outside the recipient's scope. The second prevents a previously authorized report from being served after permission changes. The stored digest binds output to snapshot contents but is not itself an access token. Guessing a job identifier should not grant access; the API must verify the user and incident scope for every fetch. The pilot's finite scorer does not exercise browser sessions or cross-organization access and therefore cannot validate those controls; they belong to the application integration tests.

A lifecycle test should include an in-flight job. If permissions are revoked after generation starts but before the worker stores the output, the worker must compare its snapshot digest with the current authorized snapshot before publishing. If the job was completed before revocation, the read path must still deny or mark it stale. There are two race windows, and closing only one is insufficient. The service's current worker and HTTP read paths both check for a digest mismatch; a future formal analysis could state linearization points and prove the absence of a served stale report under concurrent updates.

Authentication and authorization are separate. A valid browser session proves the identity of a user, while incident-scoped permissions determine which snapshots that user may read. A second-laptop model endpoint has a machine credential and TLS channel but no independent right to fetch arbitrary incidents. It receives only the snapshot the API already approved. If that machine is compromised, it can read the snapshots sent to it during its authorized operation. This is a residual risk of moving inference off the primary PC, even though no external cloud provider is contacted.

An authorization check can fail because the database is unavailable. Serving an old report in that state would convert a temporary outage into a disclosure decision. The safer response is a visible unavailable status or a fallback built only after current authorization can be established. The distinction is operationally costly: an analyst may lose an explanation during a database incident. For sensitive cross-organization data, that cost is preferable to silently reusing a stale authorization decision. The evaluation does not quantify that availability trade-off.

\section{Contract Invariants and Their Preconditions}
The reference invariant is that every displayed evidence-value sentence names an evidence item in the recipient's current snapshot with the same version and value. This follows from computing allowed tuples from that snapshot and rendering only admitted tuples. It assumes the snapshot constructor has already enforced access and that the renderer does not append unchecked model prose. The invariant does not make the sensor accurate. It does not apply to an image of an earlier report taken before a correction.

The kind invariant is that a partial-coverage field cannot be rendered as a complete observation. The expected-claim constructor maps partial coverage to an unknown claim. The template text says that unavailable evidence cannot establish absence of threat. A model that tries to convert partial coverage into a reassuring observation fails tuple membership. The strength of this rule depends on coverage being accurately represented upstream. If the incident API incorrectly labels a partial snapshot complete, the explanation layer cannot discover the omission.

The version invariant is that a completed report is usable only while its digest matches a newly authorized snapshot for the requesting user. This is enforced on retrieval and also checked by the worker before storing a result. A digest mismatch makes the job stale. The digest is an equality token over snapshot state, not a cryptographic proof of source authenticity. If two distinct security states produce an identical digest through a construction bug, this invariant fails. Testing must therefore include version changes, deactivation, and permission changes and verify that each changes the authorized view or blocks access.

The bounded-attempt invariant is that a model request is attempted at most twice: one initial generation and one repair. This prevents an unbounded autonomous loop. It says nothing about how long either call takes until a timeout is configured and enforced. An exception also ends in deterministic fallback, which keeps detection independent of model availability. The pilot's median generation time shows that the operational limit must be expressed in seconds and resource reservations, not just attempt count.

Finally, the no-action invariant is that the explanation component has no remediation tool. A suggestion to review local logs is represented as advice and rendered as a suggestion to an authorized analyst. It cannot launch a query or change an incident. This reduces the damage from a malicious instruction in evidence. It does not prevent a human analyst from acting on misleading advice, so user-facing labels and review procedures remain part of the larger safety case.

These invariants are deliberately narrow. None claims that the output is a complete, unbiased, or strategically useful account of an incident. A new predicate, renderer, or external delivery channel changes the proof obligations. Treating the verifier as a stable interface requires contract-versioning and regression tests whenever the vocabulary expands. The present study is a reproducible baseline for that work, not a claim of final coverage.
